\subsection{Regresión Ridge}
\label{subsubsec:eda_reg_rid}

El método de regresión Ridge permite detectar la multicolinealidad dentro de un modelo de regresión lineal. Fue propuesto por Hoerl y Kennard y es usado para trabajar con modelos que presentan sesgo ~\cite{dialnet5116534, ibmRidgeRegression}. 

Este método consiste en que dado una matriz $(X^T X)$ altamente condicionada o cercana a singular, esto es, que es extremadamente sensible a pequeñas variaciones de los datos de entrada, es posible agregar constantes positivas a los elementos de la diagonal para asegurar que la matriz resultante no sea altamente condicionada.

En relación con el machine learning, la regresión Ridge sirve para corregir el sobreajuste de los datos de entrenamiento. 

Cabe aclarar que la regularización es un método estadístico para reducir los errores causados por el sobreajuste de los datos de entrenamiento mientra que la regresión corrige específicamente la multicolinealidad en el análisis de regresión.

Este método resulta útil se desarrollan modelos con un gran número de parámetros.

La multicolinealidad denota cuando dos o más predictores tienen una relación casi lineal.

La regresión Ridge funciona de la siguiente manera:

\begin{itemize}
 \item Coeficientes: utiliza un estimador estándar de coeficientes matriciales por mínimos cuadrados ordinarios (MCO): $B = (X^T X)^-1 X^t y$. Esto se consigue encontrando la línea que mejor se ajuste a un conjunto de datos determinado mediante el cálculo de los coeficientes de cada variable independiente que, en conjunto, den como resultado la menor suma residual de cuadrados.
 
 \item Suma residual de cuadrados (RSS): mide la adecuación de un modelo de regresión lineal a los datos de entrenamiento. $RSS = \sum_{i=1}^{n} (Y_i - \hat{Y}_l)^2$. Esta fórmula mide la precisión de la predicción del modelo para los valores reales de los datos de entrenamiento. Si RSS = 0, el modelo predice perfectamente las variables dependientes.
\end{itemize}

En un modelo de regresión lineal, si dos o mas variables comparten una correlación lineal alta, el MCO puede devolver coeficientes erroneos siendo esto prueba de sobreajuste. 

La regresión Ridge lo que hace es modificar el MCO calculando coeficientes que tienen en cuenta predictores potencialmente correlacionados corrigiendo los coeficientes de alto valor introduciendo un término de regularización en la función RSS.

Término de regularización $L2 = ||B||^2 = B_1^2 + B_2^2 + ... + B_n^2$. Esto hace que la formula del RSS quede de la siguiente manera:
$RSS_L2 = sum_{i=1}^{n} (Y_i - \hat{Y}_l)^2 + \lambda sum_{j=1}^{p} B_j^2$.

La agregación del L2 al RSS contrarresta los coeficientes especialmente altos al reducir todos los valores de los coeficientes.